\section{Introduction}
This paper is a study of the confluence of solutions of the generic Painlev\'e-III equation to solutions of its parameter-free degeneration. The six Painlev\'e equations and their solutions, often referred to as \emph{Painlev\'e transcendents}, have been the subject of intense study. This is largely motivated by the fact that Painlev\'e transcendents are generically transcendental, and yet appear in various applications in integrable systems, integrable probability, and random matrix theory to name a few.

\subsection{B\"acklund transformations and rational solutions of Painlev\'e-III}
All Painlev\'e equations but the first are actually families of differential equations indexed by complex parameters appearing as coefficients. However, certain solutions corresponding to different parameters can be related via \emph{B\"acklund transformations}. For example, consider our main object of study, the generic Painlev\'e III equation, known as PIII($D_6$):
\begin{equation}
\dod[2]{u}{x}=\dfrac{1}{u}\left( \dod{u}{x} \right)^2-\dfrac{1}{x} \dod{u}{x} + \dfrac{\alpha u^2 + \beta}{x}+4u^3-\frac{4}{u}, \qquad \alpha,\beta \in \C.
\label{eq:PIII-$D_6$}
\end{equation}
In \cite{Gromak}, Gromak discovered that the transformation
\begin{equation}
 u(x) \mapsto \hat{u}(x) := \dfrac{2xu'(x) + 4xu(x)^2 + 4x -\beta u(x) - 2u(x)}{u(x) ( 2xu'(x) + 4x u(x)^2 +4x + \alpha u(x) + 2u(x) )}
\label{eq:Gromak-transformation}
\end{equation}
mapped solutions of \eqref{eq:PIII-$D_6$} with parameters $(\alpha, \beta)$ to solutions of \eqref{eq:PIII-$D_6$} with parameters $(\alpha + 4,\allowbreak \beta+4)$. With this one can construct from a given seed solution $(u_0, \alpha, \beta)$ a family of solutions $(u_n, \alpha + 4n, \beta + 4n)$ by iterating transformation \eqref{eq:Gromak-transformation}. The paper \cite{MCB} contains a survey of families of solutions of \eqref{eq:PIII-$D_6$} constructed using this and other B\"acklund transformations. A notable family of solutions constructed in this manner is a sequence of rational solutions $u=u_n(x;m)$ obtained from the seed function $u_0(x) \equiv 1$ and parameters $\alpha = -\beta = 4m$. This family of solutions has been numerically and analytically explored in \cite{BMS18}, and many conjectures were formulated there. While some of these were later resolved in the sequel \cite{BM20}, some conjectures remained open, see \cite[Conjectures 4 and 5]{BMS18}. Conjecture 5 is concerned with the behavior of $u_n(x;m)$ near the singular point $x = 0$. As was done in \cite{BMS18}, writing \[z = nx,\qquad U_n(z;m):= u_n(x;m)\] and considering large $n$ for fixed $m$ yields the differential equation
\[
\dod[2]{U_n}{z} = \dfrac{1}{U_n} \left( \dod{U_n}{z}\right)^2 - \dfrac{1}{z} \dod{U_n}{z} + \dfrac{4U_n^2 + 4}{z} + \mathcal{O}\big(n^{-1}\big).
\]
Formally taking the limit and denoting the limiting function $U(z;m)$ yields the parameter-free Painlev\'e-III equation, referred to as PIII($D_8$),
\begin{equation}
\dod[2]{U}{z}=\dfrac{1}{U}\left( \dod{U}{z}\right)^2-\dfrac{1}{z} \dod{U}{z} + \dfrac{4U^2 + 4}{z}.
\label{eq:PIII-$D_8$}
\end{equation}
The content of Conjecture 5 is that this convergence holds at the level of solutions, not just equations.

\subsection{Results}
To begin with, we prove Conjecture 5 from \cite{BMS18} in this work; to be more precise we establish the following theorem.
\begin{Theorem}\label{thm:limit-of-rational-solutions}
Fix $m \in \C \setminus \big(\Z + \frac{1}{2}\big)$ and let $u_n(x;m)$ be the family of rational solutions described above. There exists a unique solution $U(z)=U(z;m)$ of the Painlev\'e-{\rm III}$(D_8)$ equation \eqref{eq:PIII-$D_8$} analytic at the origin with $U(0;m)=\tan\big(\frac{\pi}{2}\big(m+\frac{1}{2}\big)\big)\neq 0$ such that
\begin{gather}
\lim_{j \to \infty} u_{2j}\left(\frac{z}{2j}; m\right) = U(z;m),\nonumber\\
\lim_{j \to \infty} u_{2j+1}\left(\frac{z}{2j+1}; m\right) = -1/U(z;m)\label{eq:even-odd-limit}
\end{gather}
for $z\notin \Sigma(m)$, where $\Sigma(m)$ denotes the union of all poles and zeros of $z\mapsto U(z;m)$. The convergence is uniform on compact subsets of $\C \setminus \Sigma(m)$.
\end{Theorem}
\begin{figure}[t]
 \centering
\includegraphics[scale=0.8]{Figures/rationaln10m1over4.png}\qquad
\includegraphics[scale=0.8]{Figures/limiting_rational.png}
 \caption{Left: the rational solution $u_{10}(x;0.25)$. Right: the limiting solution $U(z;0.25)$, where we recall the notation $z=nx$ for $n=10$. All poles of $u_{10}(x;0.25)$ are simple with residue $\frac{1}{2}$/$-\frac{1}{2}$, indicated in the plot with red/yellow circles. Likewise, all zeros of $u_{10}(x;0.25)$ are simple with derivative $2$/$-2$, indicated in the plot with pink/green squares. On the other hand, all poles and zeros of $U(z;0.25)$ have multiplicity $2$ and are marked with red circles and green squares respectively. }
 \label{fig:rational}
\end{figure}
We illustrate this theorem in Figure \ref{fig:rational}. The pictures are made using the code from \cite{FFW}, which was generously provided by the authors.

In Section \ref{sec:Frobenius-method}, we study the Maclaurin series solutions of \eqref{eq:PIII-$D_6$}; this characterizes the limiting solution of \eqref{eq:PIII-$D_8$} via its initial conditions and produces a local version of Theorem \ref{thm:limit-of-rational-solutions}, see Theorem \ref{thm:$D_6$-$D_8$-series} and Corollary \ref{cor:local-limit-of-rational-solutions} below.

The rational solutions $u_n(x;m)$ are related to the so-called Umemura polynomials $s_n(x;m)$ by the formula
\begin{equation}
u_n(x;m)=\dfrac{s_n(x;m-1)s_{n-1}(x;m)}{s_n(x;m)s_{n-1}(x;m-1)}.
\label{eq:un-polys}
\end{equation}
Indeed, a sequence of rational functions $x\mapsto s_n(x;m)$ is determined by the recurrence relation
\begin{gather}
s_{n+1}(x;m)\nonumber\\
\quad=\dfrac{(4x+2m+1)s_n(x;m)^2-s_n'(x;m)s_n(x;m)-x\big(s_n''(x;m)s_n(x;m)-s_n'(x;m)^2\big)}{2s_{n-1}(x;m)}
\label{eq:umemura-recurrence}
\end{gather}
with initial conditions
\begin{equation}\label{eq:umemura-recurrence_init}
s_{-1}(x;m)=1,\qquad s_{0}(x;m)=1.
\end{equation}
It was shown in \cite{clarkson2018constructive,Umemura} that the rational functions $s_n(x;m)$ are actually polynomials. In Section~\ref{sec:Umemura}, we use Corollary \ref{cor:local-limit-of-rational-solutions} to deduce asymptotics of the Umemura polynomials themselves. To formulate our result, we need to introduce a certain Fredholm determinant; more precisely, let $K_r\colon L^2[0, r] \to L^2[0, r]$ denote the integral operator with the continuous Bessel kernel
 	\begin{align*}
 	&K(x,y)=\frac{\sqrt{x}J_1\big( \sqrt x\big) J_0\big(\sqrt y\big)-
 	 J_0\big(\sqrt x\big) \sqrt{y}J_1\big( \sqrt y\big)}{2( x-y)}.
 	\end{align*}
For any $\lambda\in\mathbb C$, let $D_\lambda( r)$ be the Fredholm determinant
\[%\label{eq:bessel-determinant}
 	D_\lambda( r) :=\det ( \mathbf 1-\lambda K_r).
\]
It is well known (see, e.g., \cite[Chapter 24]{Lax}) that the Fredholm determinant $D_\lambda( r)$ is an entire function of $\lambda$. Since $K_r$ is a trace-class integral operator, one of several equivalent ways to define~$D_\lambda\lb r\rb$ is via the Plemelj--Smithies formula
 \begin{gather}
 D_\lambda( r)=\exp\biggl(-\sum_{\ell=1}^\infty \Tr K_r^{\ell} \frac{\lambda^\ell}{\ell}\biggr).
 \label{eq:plemelj-smithies}
 \end{gather}
The traces in \eqref{eq:plemelj-smithies} have explicit expressions as iterated integrals
\[
\Tr K_r^\ell := \int_0^r K^{(\ell)}(t, t) \dd t,
\]
where
\[
K^{(1)}(x, y) = K(x, y) \qquad \text{and} \qquad K^{(\ell)}(x, y) = \int_0^r K(x, t) K^{(\ell - 1)}(t, y) \dd t.
\]
By re-scaling the integrals to bring the $r$-dependence to the integrand and observing that $J_0 \big( \sqrt{xy} \big)$ and $\sqrt{xy}J_1\big(\sqrt{xy}\big)$ are both entire functions with respect to both $x$ and $y$, we see that $\Tr K_r^{\ell}$ and~$D_\lambda ( r)$ can be extended to analytic functions of $r$ in a neighborhood of $r=0$, and in fact~$\Tr K_r^{\ell}=\mathcal{O}\big(r^\ell\big)$ as $r\to 0$, from which we obtain $D_\lambda(0)=1$. We are now ready to state our second theorem.
\begin{Theorem}\label{thm:umemura_asym}
Fix $m \in \C \setminus \big(\Z + \frac{1}{2}\big)$. Then, there exists a small enough neighborhood of the origin, $\mathcal{G}$, such that the Umemura polynomials admit the following limits along the even and odd subsequences:
\begin{equation}
 \lim_{j\to\infty}\frac{s_{2j}\big( \frac{z}{2j+1};m \big)}{s_{2j}(0;m)}=\ee^{2\ii z}\left(\frac{U(z;m)}{U(0;m)}\right)^{-\frac{1}{4}} \sqrt{D_{\lambda(m)}\lb 32\ii z\rb},
\label{eq:umemura-asymtotics-1-bessel}
\end{equation}
and
\begin{equation}
 \lim_{j\to\infty}\frac{s_{2j-1} \big(\frac{z}{2j} ;m\big)}{s_{2j-1}(0;m)}=\ee^{2\ii z}\left(\frac{U(z;m)}{U(0;m)}\right)^{\frac{1}{4}} \sqrt{D_{\lambda(m)}\lb 32\ii z\rb},
\label{eq:umemura-asymtotics-2-bessel}
\end{equation}
where $\lambda(m)=1/\big(1+\ee^{2\pi \ii m}\big)$, the square root and fractional powers denote the principal branches taking the value $1$ at $z=0$, and the convergence is uniform for $z\in\mathcal{G}$. Furthermore,
the values of the Umemura polynomials at the origin have the leading asymptotics
\begin{gather}
 s_{2j}(0;m)\sim\sqrt{2\pi}\ee^{4\zeta'(-1)}\frac{j^{2j^2+j+\frac{m^2}{2}+\frac{m}{2}+\frac{1}{24}}\ee^{-3j^2-j}2^{2j^2+2j}(-\cos(\pi m))^j}{G\big(\frac{5}{4}+\frac{m}{2}\big)G\big(\frac{5}{4}-\frac{m}{2}\big)G\big(\frac{7}{4}+\frac{m}{2}\big)G\big(\frac{3}{4}-\frac{m}{2}\big)},\qquad j\to\infty,\label{eq:umemura-zero-asym-even}
\\
 s_{2j-1}(0;m)\sim\frac{\ee^{4\zeta'(-1)}}{\sqrt{2\pi}}\frac{j^{2j^2-j+\frac{m^2}{2}+\frac{m}{2}+\frac{1}{24}}\ee^{-3j^2+j}2^{2j^2}(\cos(\pi m))^j}{G\big(\frac{3}{4}+\frac{m}{2}\big)G\big(\frac{3}{4}-\frac{m}{2}\big)G\big(\frac{5}{4}+\frac{m}{2}\big)G\big(\frac{1}{4}-\frac{m}{2}\big)},\qquad j\to\infty,\label{eq:umemura-zero-asym-odd}
\end{gather}
in which $G$ denotes the Barnes $G$-function and $\zeta$ denotes the Riemann zeta function.
\end{Theorem}
In fact, one can check that the expressions on the right-hand side of \eqref{eq:umemura-asymtotics-1-bessel} and \eqref{eq:umemura-asymtotics-2-bessel} admit analytic continuation from a neighborhood of $z=0$ to the whole $z$-plane. Although it does not follow from our proof, this suggests that the neighborhood $\mathcal{G}$ can be taken to be an arbitrary bounded set.

Our analysis of series solutions in Section \ref{sec:Frobenius-method} points to a more general statement about the coalescence of solutions of \eqref{eq:PIII-$D_6$} to solutions of \eqref{eq:PIII-$D_8$}. The technical result leading to Theorem~\ref{thm:limit-of-rational-solutions} by Maclaurin series (see Theorem \ref{thm:$D_6$-$D_8$-series} below) applies not only to rational solutions, but to all sequences of solutions with initial conditions converging to finite, nonvanishing limits. This, however, is a serious limitation since $x = 0$ is a singular point of Painlev\'e-III, and generic solutions of \eqref{eq:PIII-$D_6$} will be singular at this point and behave like $u(x)\simeq ax^p$, $|{\re}(p)|< 1$. More specifically, based on symbolic computation we expect the asymptotic expansion for solutions of \eqref{eq:PIII-$D_6$} in the form
\[
u(x)\sim \sum_{k=0}^{\infty}\sum_{l=0}^{k+1}\big(b_{kl}x^{2k+(2l-1)p}+c_{kl}x^{2k+1+2lp}\big) \qquad \text{as}\quad x\to 0.
\]
To tackle this issue, we develop a second approach that avoids series expansions and instead relies on the isomonodromy representation of the Painlev\'e transcendents. It was first discovered by Garnier \cite{MR1509146} and further explicated by Jimbo and Miwa in \cite{MR625446} that Painlev\'e equations can be formulated as monodromy-preserving, or isomonodromic, deformations of corresponding~${2 \times 2}$ first-order systems of differential equations. This allows one to characterize solutions of a~given Painlev\'e equation in terms of a $2\times 2$ Riemann--Hilbert problem. Such a monodromy representation was obtained for rational solutions of Painlev\'e-III($D_6$) in \cite{BMS18}. From this point of view, one can show that for fixed $\alpha, \beta \in \C$, the solutions of \eqref{eq:PIII-$D_6$} are parametrized by triples~$(x_1, x_2, x_3) \in \C^3$ on the cubic surface, known as the \emph{monodromy manifold}, given by
\begin{equation}
 \label{eq:cubic-D6} x_1x_2x_3+x_1^2+x_2^2+x_2 \bigl(\ee^{-{\ii\pi \alpha}/{4}}-\ee^{-{\ii\pi \beta}/{4}} \bigr)+x_1 \bigl(1-\ee^{-{\ii\pi(\alpha+ \beta)}/{4}}\bigr)-\ee^{-{\ii\pi(\alpha+ \beta)}/{4}}=0.
\end{equation}
The exponential constants appearing as coefficients in \eqref{eq:cubic-D6} will appear in multiple equations, making it convenient to introduce the notation
\begin{equation}
 e_0:=\ee^{{\ii\pi\alpha}/{8}}\neq 0 \qquad \text{and} \qquad e_\infty:=\ii\ee^{-{\ii\pi\beta}/{8}}\neq 0.
 \label{eq:e0-einfty-alpha-beta}
\end{equation}
In Section \ref{sec:monodromy-rep-$D_6$}, we reproduce the derivation of the cubic surface \eqref{eq:cubic-D6} carried out in \cite{PS} and connect the quantities $x_i$ with other invariant quantities that appear in the Riemann--Hilbert Problem~\ref{rhp:initial} associated with PIII($D_6$). In Section \ref{sec:monodromy-rep-$D_8$}, we present an analogous parametrization of solutions of the $D_8$ degeneration \eqref{eq:PIII-$D_8$} of PIII in terms of triples $(y_1, y_2, y_3) \in \C^3$ appearing in the Riemann--Hilbert Problem~\ref{rhp:D8} and satisfying
\begin{equation}
\label{eq:cubic-D8}
y_1 y_2 y_3 + y_1^2 + y_2^2 + 1 = 0.
\end{equation}
Away from its singular points, we parametrize points $(x_1,x_2,x_3)$ on the cubic surface \eqref{eq:cubic-D6} using parameters $e_1$, $e_2$ appearing naturally from the point of view of the Riemann--Hilbert problem. In fact, $e_1^2$, $e_1^{-2}$ are eigenvalues of a certain monodromy matrix for a circuit about the origin for a linear system, see \eqref{eq:generic-system}. The parameter $e_2$ appears in the connection matrix for the same system, see \eqref{eq:C-zero-infty-diagonalization}.
We call $(e_1,e_2)$ monodromy parameters.
\begin{Definition}[see Section \ref{sec:monodromy-rep-$D_6$} for details]\label{def:generic}\samepage
We say the monodromy parameters $(e_1, e_2)$ are \emph{generic} if
 \begin{enumerate}\itemsep=0pt
 \item[(i)] $e_1^4 \neq 1$,
 \item[(ii)] $e_1 e_2 \neq 0$,
 \item[(iii)] $e_1^2 \neq e_\infty^{\pm 2}$ and $e_1^2 \neq e_0^{\pm 2}$.
 \end{enumerate}
\end{Definition}
Before moving on, we pause to make a few observations.
\begin{itemize}\itemsep=0pt
 \item Condition (ii) implies that generic monodromy parameters are nonvanishing, hence we may write
 \begin{gather}
 e_1=\ee^{\ii\pi\mu}\qquad \text{and}\qquad e_2=\ee^{\ii\pi\eta}. \label{eq:e1-e2-mu-eta}
 \end{gather}
 \item Since $e_1^2$, $e_1^{-2}$ are the essential quantities related to the complex parameter $e_1$ in our parametrization, and the former are insensitive to a change in sign of $e_1$, we may take~$e_1$ to be in the right half plane; in view of \eqref{eq:e1-e2-mu-eta}, this corresponds to $-\frac12 < \re(\mu) \leq \frac12$. Note that the choice of including the upper versus the lower endpoint of this range is arbitrary.
 \item Due to $e_1^2$, $e_1^{-2}$ being eigenvalues of the same matrix (see \eqref{eq:Stokes-products-eigenvectors} below), the parameters~$\mu$,~$-\mu$ correspond to the same solution of \eqref{eq:PIII-$D_6$}. While we could restrict to parameters where $\re(\mu) >0$ (say), we choose not to and in turn arrive at slightly simpler formul\ae; see Remark \ref{remark:e2-mu-minus-mu-change} below.
 \item It turns out that the parameter $e_2$ is determined up to a sign as well, see Remark \ref{remark:change-sign-e2} below. As such, we take it to be in the right half plane as well, or $-\frac12 < \re(\eta) \leq \frac12$.
\end{itemize}
With this in mind, we can now state a more general theorem.
\begin{Theorem}\label{thm:general}
 Let $u_0$ be the solution of \eqref{eq:PIII-$D_6$} corresponding to monodromy data $(\alpha, \beta, x_1, x_2, x_3)$ parametrized by generic monodromy parameters $(e_1,e_2)$ using formul\ae\ $($consistent with \eqref{eq:cubic-D6}$)$
 \begin{gather}%\label{eq:x2}
 x_1=
 \frac{e_1^2 \big(e_0^2 e_2^2e_\infty^2 \big(e_1^2 e_\infty^2-1\big)+e_0^2 e_1^2-1\big)\big(\big(e_0^2 e_1^2-1\big)^2 + e_0^2e_2^2e_\infty^2 \big(e_0^2-e_1^2\big) \big(e_\infty^2-e_1^2\big)\big)}{e_0^4 e_2^2e_\infty^2\big(1-e_0^2 e_1^2\big)\big(e_1^4-1\big)^2 },\label{eq:x1}
\\
 x_2= \frac{\big(e_0^2 e_2^2 e_1^2e_\infty^2 \big(e_1^2-e_\infty^2\big)+1-e_0^2 e_1^2\big) \big(\big(e_0^2 e_1^2-1\big)^2\!+e_0^2 e_1^2e_2^2 e_\infty^2 \big(e_0^2-e_1^2\big) \big(e_1^2 e_\infty^2\!-1\big)\big)}{e_0^4e_2^2e_\infty^2 \big(1-e_0^2 e_1^2\big) \big(e_1^4-1\big)^2 }, \!\!\!\!
\\
 x_3=
 e_1^2+\frac{1}{e_1^2},\label{eq:x3}
\end{gather}
and let $U(z)=U(z;y_1,y_2,y_3)$ denote the solution of \eqref{eq:PIII-$D_8$} with monodromy data $($consistent with \eqref{eq:cubic-D8}$)$
 \begin{gather}
 y_1 = \ii\sqrt{\frac{e_\infty^2-e_1^2}{1-e_0^2 e_1^2}} \cdot\frac{ \big(1-e_0^2 e_1^2+e_0^2e_1^6e_2^2e_
 \infty^2\big(e_1^2-e_\infty^2\big)\big)}{e_0 e_1 e_2 e_\infty\big( e_\infty^2-e_1^2\big) \big(e_1^4-1\big) },
 \label{eq:y1}
\\
 y_2 = \ii\sqrt{\frac{e_\infty^2-e_1^2}{1-e_0^2 e_1^2}} \cdot\frac{e_1 \big(1-e_0^2 e_1^2+e_0^2e_1^2e_2^2e_
 \infty^2\big(e_1^2-e_\infty^2\big)\big)}{e_0e_2 e_\infty \big( e_\infty^2-e_1^2\big) \big(e_1^4-1\big)},
 \label{eq:y2}
 \\
 y_3 = -e_1^2-\frac{1}{e_1^2}.
 \label{eq:y3}
 \end{gather}
If $u_n$ is the $n$th iterate of $u_0$ under transformation \eqref{eq:Gromak-transformation}, then for $z\notin \Sigma(y_1,y_2,y_3)$
 \begin{gather*}
 \lim_{j \to \infty} u_{2j}(z/2j) = U(z;y_1,y_2,y_3),\\ \lim_{j \to \infty} u_{2j+1}(z/(2j+1)) = -1/U(z;y_1,y_2,y_3),
 \end{gather*}
where convergence is uniform on compact subsets of $\C\setminus \Sigma(y_1, y_2, y_3)$ slit along $\arg(z)= \pm \pi$ and~$\Sigma(y_1,y_2,y_3)$ is the union of all poles and zeros of $z\mapsto U(z;y_1,y_2,y_3)$.
\end{Theorem}
There is nothing fundamental about the exclusion of $\arg(z) = \pm \pi$; in fact, the Riemann--Hilbert analysis below can be continued onto the universal cover of $\C \setminus \{0\}$ with a suitable extension of the set $\Sigma(y_1, y_2, y_3)$. Similar observations apply to Proposition \ref{prop:un-zero-asymptotics} and Theorem~\ref{thm:D8-asymptotics-zero} below.

We illustrate Theorem \ref{thm:general} for solutions that are not single-valued near the origin in Figure \ref{fig:random}. Note that while the point $(x_1,x_2,x_3)\in\mathbb{C}^3$ on the monodromy manifold \eqref{eq:cubic-D6} only depends on the squares of $e_1$, $e_2$, the point $(y_1,y_2,y_3)\in\mathbb{C}^3$ on the monodromy manifold \eqref{eq:cubic-D8} of the limiting solution $U(z;y_1,y_2,y_3)$ of \eqref{eq:PIII-$D_8$} has a sign ambiguity in the coordinates $y_1$ and $y_2$. However, if either $e_1$ or $e_2$ changes sign, then the signs of $y_1$ and $y_2$ change together, and it turns out that the triples $(y_1,y_2,y_3)$ and $(-y_1,-y_2,y_3)$ both lie on the surface \eqref{eq:cubic-D8} together and correspond to the same solution of \eqref{eq:PIII-$D_8$}; see Remark~\ref{rem:minus-y1-y2} below. Similarly, there is no need for us to specify the sign of the square roots in \eqref{eq:y1}--\eqref{eq:y2} provided they are both taken to be the same. One might expect a similar ambiguity to arise from the replacement of $e_1^2 \mapsto e_1^{-2}$, since both are eigenvalues of the same matrix, but it turns out that $(x_1, x_2, x_3)$ is invariant under this change provided~$e_2$ is appropriately modified, and $(y_1, y_2, y_3)$ remains invariant up to the sign ambiguity described above, see Remark~\ref{remark:e2-mu-minus-mu-change} below.

\begin{figure}[t]
 \centering
\includegraphics[scale=0.8]{Figures/randomd6.png}\qquad
\includegraphics[scale=0.8]{Figures/randomd8.png}
 \caption{Left: the solution $u_{10}(x)$ of \eqref{eq:PIII-$D_6$} generated by ten iterations of \eqref{eq:Gromak-transformation} with seed $u_0(x)$ corresponding to monodromy data $\mu=0.23+0.39\ii$ (see \eqref{eq:e1-e2-mu-eta}), $e_2=-0.45-0.96\ii$ and $\alpha=40.5+0.63\ii$, $\beta=40.98+0.59\ii$. Right: the limiting solution $U(z)$ of \eqref{eq:PIII-$D_8$}. The labeling of poles and zeros is the same as in Figure \ref{fig:rational}. Note that both $u_{10}(x)$ and $U(z)$ are branched at the origin.}
 \label{fig:random}
\end{figure}

The proof of Theorem \ref{thm:general} is given in Section \ref{sec:proof}, and relies on Riemann--Hilbert analysis. The idea of the proof is to use parametrices constructed out of confluent hypergeometric functions near zero and infinity to reduce the setup to a Riemann--Hilbert Problem~\ref{rhp:Q} on the circle. After some additional transformations, the problem allows taking a large $n$ limit which gives us a~Riemann--Hilbert Problem~\ref{rhp:Rhat-even/odd} with a~jump on the circle in terms of Bessel functions. Further transformations using parametrices constructed out of Bessel functions simplify the jump and we arrive at a Riemann--Hilbert Problem~\ref{rhp:D8} for Painlev\'e-III$(D_8)$. In Section~\ref{sec:suleimanov-solution-connection}, we transform this into another Riemann--Hilbert Problem~\ref{rhp:S-hat} for \eqref{eq:PIII-$D_8$} already known in the literature. It is worth pointing out that even in the case of rational solutions, the Painlev\'e-III$(D_6)$ Riemann--Hilbert Problem~\ref{rhp:initial} exhibits Stokes phenomenon near both singular points \cite{BMS18} and hence requires the use of confluent hypergeometric parametrices to desingularize the problem before passing to the limit.

While the formul\ae\ for $y_i$ are daunting, they drastically simplify in the case of the rational solutions, where $u_0$ has monodromy data parametrized by
\begin{gather}
 \alpha = -\beta = 4m, \qquad e_0^2 = -e_\infty^2 = \ee^{\ii \pi m},\qquad
 e_1^2 = \ii, \qquad
 e_2 = \sqrt{\ee^{-2\pi \ii m}\dfrac{1 - \ii \ee^{\pi \ii m}}{1 + \ii \ee^{\pi \ii m}}},
\label{eq:monodromy-data-rational}
\end{gather}
see Section \ref{sec:rational-solutions-parameters} for details. With parameters chosen as in \eqref{eq:monodromy-data-rational}, the genericity conditions in Definition \ref{def:generic} imply $m \in \C \setminus \big(\Z + \frac{1}{2}\big)$. Then, we have
\[
y_1 = \frac{\ii \ee^{\ii\pi m}}{\sqrt{1+ \ee^{2\pi \ii m}}}, \qquad y_2 = \frac{\ii }{\sqrt{1+ \ee^{2\pi \ii m}}}, \qquad y_3 = 0.
\]
One can check that with these choices of $\alpha$, $\beta$, $e_1,$ and $e_2$ we have $U(z;y_1,y_2,y_3) = U(z;m)$, cf.\ Theorems \ref{thm:limit-of-rational-solutions} and \ref{thm:general}.

By further specializing $U(z;m)$ to $m \in \ii\R+ \Z$, we arrive at highly symmetric solutions of~PIII($D_8$) which have appeared in various works in nonlinear optics \cite{suleimanov} and as a limiting object of various families of solutions to the focusing nonlinear Schr\"odinger equation in different regimes \cite{MR4007631, BLM20}. Furthermore, these solutions can be identified with pure imaginary solutions of the radial reduction of sine-Gordon equation, see, e.g., \cite[Chapter 13]{FIKN}. It is interesting that they are related to another limiting object appearing in the random matrix theory -- the Bessel kernel determinant. The explicit relation is described in Corollary \ref{cor:bessel-relation} below.

A consequence of the analysis in Section \ref{sec:proof} below is a description of the behavior near the origin of solutions $u(x)$ of \eqref{eq:PIII-$D_6$} corresponding to generic monodromy parameters $(e_1, e_2)$.
\begin{prop}
\label{prop:un-zero-asymptotics}
 Let $u(x)$ be the solution of Painlev\'e-{\rm III}$(D_6)$ equation associated to $\mu, \eta \in \C$ via the generic monodromy parameters given in \eqref{eq:e1-e2-mu-eta} with $-\frac{1}{2}<\re(\eta)\leq\frac{1}{2}$. If $0<|{\re}(\mu)|<\frac{1}{2}$, then it holds that
 \begin{align}
 u(x) ={}& -\frac{ \Gamma (1 - 2\epsilon \mu)^2 \Gamma \left( \epsilon\mu -\frac{\alpha}{8}\right) \Gamma \big( \epsilon \mu+\frac{\beta}{8} + \frac{1}{2} \big)}{\Gamma (2 \epsilon \mu )^2 \Gamma \bigl(-\epsilon \mu-\frac{\alpha}{8}+1\bigr) \Gamma \bigl( - \epsilon \mu +\frac{\beta}{8} + \frac{1}{2} \bigr)} \nonumber\\
 &\times \Bigg(\dfrac{e_{0}^2 e_2^2 e_{\infty}^2\big(e_{0}^2-e_{1}^2\big) \big(e_{1}^2-e_{\infty}^2\big)}{\big(e_{0}^2 e_{1}^2-1\big)^2}\Bigg)^\epsilon
 x^{4 \epsilon \mu - 1} \big(1 + \mathcal{O}(x^\delta) \big)
 \label{eq:u-0-leading}
 \end{align}
 as $x\to 0$ with $|{\arg}(x)|<\pi$ where $\delta=\min(1,2-4\re(\mu))$ and $\epsilon = \sgn (\re (\mu))$.
\end{prop}
%Note that the boundary cases where $\re(\mu)=\frac{1}{2}$ are not covered, and if they were included the error term would not be relatively small as one would have $\delta=0$.
Proposition \ref{prop:un-zero-asymptotics} appeared in \cite[Theorem 3.2]{Jimbo} and its derivation is given in \cite{Kitaev87} for an equivalent, degenerate Painlev\'e-V equation. We present its proof using a Riemann--Hilbert approach in Section \ref{sec:prop-1-proof}, which follows the steps of the proof of Theorem \ref{thm:general}. The case $\re(\mu) = 0$ can be handled similarly, but we exclude it here because two distinct terms arise at the same leading order resulting in a more complicated formula. From this formula one can see that if $\re(\mu)=0$ the solution can exhibit sinusoidal oscillations with frequency diverging as $x^{-1}$ consistent with an essential singularity at the origin.

To apply Proposition \ref{prop:un-zero-asymptotics} to the rational solutions \eqref{eq:un-polys}, or more generally to the sequence of B\"acklund iterates starting from any seed solution of \eqref{eq:PIII-$D_6$}, requires knowledge of the corresponding sequence of monodromy data. This is the content of the following proposition, which we prove in Section \ref{sec:schlesinger}.
\begin{prop}\label{prop:schlesinger}
Let $u_0(x)$ be the solution of \eqref{eq:PIII-$D_6$} with parameters $(\alpha,\beta)$ and monodromy data $(\mu,\eta)$ $($see \eqref{eq:e1-e2-mu-eta}$)$ with $-\frac{1}{2}<\re (\mu), \re(\eta)\le\frac{1}{2}$. Then, the B\"acklund iterates $u_n(x)$ are parametrized by the following monodromy data
\[
e^2_{1, n} = \ee^{2\pi \ii \mu_n} ,\qquad e_{2, n} = e_2, \qquad e_{0,n}=\ee^{{\ii\pi(\alpha+4n)}/{8}}, \qquad e_{\infty,n}=\ii\ee^{-{\ii\pi(\beta+4n)}/{8}}, \]
where\footnote{To ensure $-1/2<\re(\mu_n)\leq 1/2$, we set $\epsilon = -1$ in the case where $\re(\mu) = 0$.}
\[
\mu_n = \begin{cases} \mu, & n \in 2\Z, \\ \mu - \frac{\epsilon}{2}, & n + 1 \in 2\Z, \end{cases} \qquad \text{and}\qquad \epsilon = \sgn (\re (\mu)).
\]
\end{prop}

One notable application of Propositions \ref{prop:un-zero-asymptotics} and \ref{prop:schlesinger} is the case corresponding to the rational solutions of PIII($D_6$) described above. In \cite{clarkson2018constructive}, the authors found a product formula for $u_n(0;m)$ (see \eqref{eq:u-n-zero-rational} in Section \ref{sec:Frobenius-method}). Applying Propositions~\ref{prop:un-zero-asymptotics} and \ref{prop:schlesinger} to this case yields the closed-form formula
\begin{equation}\label{eq:u-n-rat-leading}
u_n(0;m)=\frac{\Gamma\big(\frac{1}{4}-\frac{m}{2}-\frac{n}{2}\big)}{\Gamma \big(\frac{1}{4}-\frac{m}{2}+\frac{n}{2} \big)}\frac{\Gamma \big(\frac{3}{4}-\frac{m}{2}+\frac{n}{2} \big)}{\Gamma \big(\frac{3}{4}-\frac{m}{2}-\frac{n}{2} \big)}.
\end{equation}
Another observation is that the expression on the right-hand side of \eqref{eq:u-0-leading} in Proposition \ref{prop:un-zero-asymptotics} evaluated at the $n$-dependent monodromy data from Proposition \ref{prop:schlesinger} and at argument $x=\frac{z}{n}$ has a finite limit along even and odd subsequences of $n$. The limiting expressions relate to the behavior of $U(z;y_1,y_2,y_3)$, which we can take from the literature:
\begin{Theorem}[\cite{FIKN,IN,N}] \label{thm:D8-asymptotics-zero}
 Let $U(z;y_1,y_2,y_3)$ be the solution of the Painlev\'e-{\rm III}$(D_8)$ equation~\eqref{eq:PIII-$D_8$} associated to $(y_1,y_2,y_3) \in \C^3$ parametrized by generic monodromy parameters $(e_1,e_2)$ using formul\ae\ \eqref{eq:y1}--\eqref{eq:y3}. Then, it holds that
 \begin{equation*}
 U(z) = -\frac{ \Gamma (1 - 2\epsilon \mu)^2 }{\Gamma (2\epsilon \mu )^2 2^{4\varepsilon\mu-1} } z^{4 \epsilon \mu - 1} (1 + \mathcal{O}(z) )\Bigg(\dfrac{e_{0}^2 e_2^2 e_{\infty}^2 \big(e_{\infty}^2-e_{1}^2\big)}{
 \big(e_{0}^2 e_{1}^2-1\big)}\Bigg)^{\epsilon}
 % \label{eq:v-leading}
 \end{equation*}
 as $z\to 0$ with $|{\arg}(z)|<\pi$.
\end{Theorem}
We pause to note that coalescence between Painlev\'e equations has long been in the literature; a coalescence diagram of all six Painlev\'e equations already appeared in Okamoto's work \cite{O1}, and was later expanded on in \cite{OO}. Later, a geometric interpretation of the coalescence was given in \cite{CMR}. That being said, the above degenerations are carried out on the level of the differential equation, so that given a solution of a Painlev\'e equation, one does not have a characterization of the solution one arrives at under the coalescence procedure. Confluence on the level of the solutions of the differential equation has also appeared in the literature; one of the most interesting examples is the merging of regular singularities and corresponding creation of an irregular singularity. This process was studied in the works \cite{Glutsuk,Kitaev06}. In the PhD thesis \cite{Horrobin} the confluence was studied in more detail in the cases Painlev\'e VI $\to$ Painlev\'e V and Painlev\'e V $\to$ Painlev\'e III$(D_6)$. In the works \cite{Kapaev94,Kapaev_Kitaev} the authors considered a transition from Painlev\'e II $\to$ Painlev\'e I that is different in nature.
